\chapter{Exercises and Solutions}

The questions appear at the end of the chapters. Solutions below show the method and a rounded result; a different result can be correct if the stated convention differs and is documented.

\section*{Chapters 1--4}
\begin{enumerate}
  \item[1.1] $(47-50+1.50)/50=-0.03$, so the return is -3 percent.
  \item[1.2] Accounting identity is a classification identity: assets equal the recorded claims on them. Market values can move without immediate accounting remeasurement.
  \item[1.3] $\E[R]=0.4(0.20)+0.6(-0.10)=0.02$. Variance is $0.4(0.18)^2+0.6(-0.12)^2=0.0144$, so volatility is 12 percent.
  \item[1.4] Funding, taxes, bid-ask spreads, short-sale constraints, margin, or counterparty default can prevent execution.
  \item[2.1] Primary market: the company sells newly issued shares and receives proceeds. Secondary market: investors trade existing shares and the company normally receives no proceeds.
  \item[2.2] It funds long-term activity with short-term claims, creating useful liquidity transformation but a mismatch if claims are withdrawn early.
  \item[2.3] Clearing can reduce bilateral settlement and replacement-cost exposure; it can concentrate default-fund, margin, operational, and liquidity risks at the central counterparty.
  \item[2.4] International standards require domestic legislation, supervision, and implementation choices.
  \item[3.1] Operating margin $=75/500=15\%$. ROE $=40/160=25\%$.
  \item[3.2] Receivables rising by 12 reduces CFO by 12; payables falling by 4 reduces CFO by 4. Total effect: CFO falls by 16.
  \item[3.3] $0.08(1.5)(2)=0.24$, or 24 percent.
  \item[3.4] Different accounting standards, lease treatment, asset age, intangibles, cyclicality, and market versus book values can change the ratio's meaning.
  \item[4.1] $10,000/1.06^4\approx7,920.94$.
  \item[4.2] Exact real return $=1.07/1.02-1\approx4.90\%$; approximation is 5 percent.
  \item[4.3] $NPV=-1000+600/1.1+600/1.1^2\approx41.32$.
  \item[4.4] The infinite growing cash-flow series converges only when $r>g$.
\end{enumerate}

\section*{Chapters 5--9}
\begin{enumerate}
  \item[5.1] $P=60/1.05+1,060/1.05^2\approx1,018.59$.
  \item[5.2] $\Delta P/P\approx-4.2(0.0025)=-1.05\%$.
  \item[5.3] Duration is linear; convexity captures the second-order curvature in the price-yield relationship.
  \item[5.4] $0.02(0.45)(5,000,000)=45,000$.
  \item[6.1] $3/(0.09-0.03)=50$.
  \item[6.2] $F=7(1.04)/(1.02)\approx7.1373$.
  \item[6.3] $(105-2)/10=10.30$ per unit.
  \item[6.4] Storage, financing, insurance, seasonality, and convenience yield can produce a positive carry even with flat expected spot.
  \item[7.1] Sunk cost: a feasibility study already paid for. Opportunity cost: rent forgone by using an owned building for the project.
  \item[7.2] No taxes and no distress costs are two assumptions. Bankruptcy costs make debt choice relevant.
  \item[7.3] $0.6(0.10)+0.4(0.05)(0.75)=0.075$, or 7.5 percent.
  \item[7.4] Staging creates the right to learn and expand only in favourable states, while avoiding some downside in unfavourable states.
  \item[8.1] $0.3(0.10)+0.7(-0.04)=0.002$, or 0.2 percent.
  \item[8.2] Uncorrelated variables can still have nonlinear dependence; independence implies zero correlation under finite moments, but the reverse is not true.
  \item[8.3] It is the model-implied change in expected excess return for a one-unit change in market excess return.
  \item[8.4] Some strategies will look significant by chance when many tests are run; the nominal 5 percent threshold no longer controls the family-wise error.
  \item[9.1] Variance $=0.5^2(0.20)^2+0.5^2(0.10)^2=0.0125$; volatility $\approx11.18\%$.
  \item[9.2] $\beta=0.018/0.0225=0.8$.
  \item[9.3] The factor model must be correctly specified and the error must be appropriately unrelated to the factor; otherwise the intercept can capture omitted risk or bias.
  \item[9.4] Expected returns and covariances are estimated with error; an unconstrained optimiser can exploit small estimation differences with extreme weights.
\end{enumerate}

\section*{Chapters 10--14}
\begin{enumerate}
  \item[10.1] Weak form concerns past market data; semi-strong form includes all public information.
  \item[10.2] Publication can lead to competition, arbitrage, changed trading costs, data-mining corrections, or a regime change.
  \item[10.3] A rule may be selected for fit, contain omitted variables, or depend on implementation assumptions; success alone does not identify a causal mechanism.
  \item[10.4] Sample period, universe, timing, benchmark, costs, turnover, capacity, and risk statistics are examples.
  \item[11.1] $Z$ is a standard-normal shock that generates the random terminal outcome.
  \item[11.2] $\mathbb Q$ is a pricing device that changes probabilities while preserving no-arbitrage pricing; it is not a claim about actual beliefs.
  \item[11.3] $0.008\sqrt{252}\approx0.1270$, or 12.70 percent.
  \item[11.4] The squared lagged shock increases the next conditional variance by $\alpha\varepsilon_{t-1}^2$.
  \item[12.1] $100e^{0.04(0.5)}\approx102.02$.
  \item[12.2] Initial margin is collateral against performance; losses can exceed it and variation margin can create further calls.
  \item[12.3] It scales the interest payments and therefore the exposure; principal is not necessarily exchanged.
  \item[12.4] Location mismatch and maturity mismatch are examples.
  \item[13.1] Call payoff is $\max(62-50,0)=12$; put payoff is zero.
  \item[13.2] Long stock plus a put replicates a call plus the present value of the strike, giving equality under no-arbitrage.
  \item[13.3] Hold half a share for each call and finance the remaining replicating position with the bond.
  \item[13.4] Real markets exhibit volatility smiles and skews; a single BSM parameter is an implied quote for each contract, not a universal constant.
  \item[14.1] Standard error falls as $1/\sqrt M$, so large accuracy gains require disproportionately more paths.
  \item[14.2] Position bounds, turnover limits, leverage caps, liquidity limits, and factor-neutrality constraints are examples.
  \item[14.3] Calibration error is mismatch between model and observed calibration targets; model error is the failure of the model structure to describe relevant reality.
  \item[14.4] Without timestamps, a researcher cannot verify that the signal used only information available at the decision time.
\end{enumerate}

\section*{Chapters 15--18}
\begin{enumerate}
  \item[15.1] VaR does not reveal the size or average of losses beyond the 99th percentile.
  \item[15.2] Market liquidity concerns trading without price impact; funding liquidity concerns meeting cash obligations.
  \item[15.3] A counterparty becomes weakest exactly when its obligation to you is largest, such as a derivative exposure rising as its credit spread widens.
  \item[15.4] A scenario preserves a coherent joint shock and can explore states not represented in a short historical sample.
  \item[16.1] Using an earnings number revised after the trading date is look-ahead bias.
  \item[16.2] A mid-price assumes immediate execution without paying the spread or moving the market; large orders cannot generally obtain it.
  \item[16.3] It is the largest percentage decline from a previous wealth peak to a later trough.
  \item[16.4] Kill switch, order-size limits, stale-data checks, reconciliation, audit logs, and incident response are examples.
  \item[17.1] Shock, margin call, forced sale, lower price/collateral, tighter funding, and further forced sale.
  \item[17.2] Insolvency means liabilities exceed asset value; illiquidity means obligations cannot be met when due, even if assets might cover them over time.
  \item[17.3] Questions include disparate error costs, explainability, appeal, privacy, data provenance, and whether users consented to the risk.
  \item[17.4] The label may use a narrow definition, omit transition risk, rely on poor data, or address impact rather than financial risk.
  \item[18.1] Build revenue and margin forecasts, calculate after-tax UFCF, discount each year, calculate terminal value, subtract net debt, and divide by shares.
  \item[18.2] Terminal value is highly sensitive when $WACC-g$ is small, so show a two-dimensional sensitivity and avoid false precision.
  \item[18.3] Review debt service, covenant headroom, refinancing, downside cash flows, and the value of flexibility.
  \item[18.4] Compare marginal contribution to portfolio return, covariance, factor exposure, liquidity, and concentration rather than standalone return alone.
\end{enumerate}

